\documentclass{article}
\usepackage{geometry}
\usepackage{graphicx} % Required for inserting images
\usepackage{amsmath, amsthm, amssymb}
\usepackage{parskip}
\newgeometry{vmargin={15mm}, hmargin={24mm,34mm}}
\theoremstyle{definition} 
\newtheorem{definition}{Definition}

\newtheorem{theorem}{Theorem}[section]
\newtheorem{lemma}[theorem]{Lemma}
\newtheorem{corollary}{Corollary}[theorem]


\title{Gathmann CA Solutions}
\date{March 2025}
\author{Boran Erol}

\begin{document}

\maketitle

\section{Exercise 0.15(a,d)}

Let $ f \in \mathbb{R}[x] $ be a polynomial vanishing on $ \mathbb{Z} $.
Then, $ f $ is the zero polynomial, so it vanishes everywhere. Then, $ \mathbb{Z} $
can't be a variety since then it would be bijective to $ \mathbb{R} $,
but the cardinalities don't match. The same argument works for part (d)
as well.

\section{Exercise 0.15(b)}

The coordinate ring of $ \mathbb{R} \backslash \{ 0 \} $ is $ \mathbb{R}[x] $,
so it's not a variety.

\section{Exercise 0.15(c)}

% TODO:

\section{Exercise 0.15(e)}

% TODO:

\section{Exercise 1.8}

Let $ x \in I_{1} \cap \cdots \cap I_{n} $.
We'd like to show that $ x \in I_{1} \cdots I_{n} $.

Let's first show it for $ n = 2 $. Since $ I,J $ are coprime,
there exists $ a_{1}, a_{2} \in R $ and $ z_{1} \in I $ and $ z_{2} \in J $
such that $ a_{1}z_{1} + a_{2}z_{2} = 1 $ by the partition of unity.

Then, $ x = x(a_{1}z_{1} + a_{2}z_{2}) = x a_{1} z_{1} + x a_{2} z_{2} $.

Notice that both terms are in the product, so we're done.

Now, assume the statement holds for $ n - 1 $. Then, we have that $ z $
is in $ I_{1} \cdots I_{n - 1} $.

Let's now show that $ I_{1} \cdots I_{n - 1} $ and $ I_{n} $ are coprime.
This is sufficient since we can apply the $ n = 2 $ case to finish the proof.

% TODO: 

\section{Exercise 1.9}

Let $ I = \langle x_{1}, \ldots, x_{n} \rangle $.

Let $ J = \langle x_{1}^{2}, \ldots, x_{n}^{2} \rangle $.

Consider $ R = K[x_{1}, \ldots, x_{n}]/J $. Since $ I $ contains $ J $,
the image of $ I $ is nonzero in $ R $. Then, a generator for $ I $ in 
$ K[x_{1}, \ldots, x_{n}] $ would be a generator for $ I $ in $ R $.
But $ I $ is now a subspace of the $ n $-dimensional vector space $ R $
over $ K $, so we need $ n $ generators.

\section{Exercise 1.11(a)}

Let $ J $ be an ideal in $ \mathbb{C}[x] $.
Since $ \mathbb{C}[x] $ is a PID, $ J = (f) $.
Let $ f = \pi_{i = 1}^{n} (x - a_{i}) $.
If $ J $ is a radical ideal, $ f $ is separable over $ \mathbb{C} $.
Thus, the vanishing ideal is simply the roots of $ f $.

Recall that $ \mathbb{A}^{1}_{\mathbb{C}} $ has the cofinite topology
as its Zariski topology. Thus, to every subvariety, we can assign
the ideal generate by the polynomial that has the subvariety as its
roots.

These two assignments are clearly inverses of each other.

\section{Exercise 1.11(b)}

The vanishing locus of $ (x^{2} + 1) $ is empty, but it's not a unit ideal.

\section{Exercise 2.11(a)}

Let $ f \in R / I_{a} $ be non-zero.

Then, $ f(a) \neq 0 $ and $ f $ is continuous.

Then, $ f $ is non-zero around a small enough neighborhood of $ a $.

\section{Exercise 2.11(b)}

Let $ f_{1}, \ldots, f_{n} \in R $ such that $ V(f_{1}, \ldots, f_{n}) = \varnothing $.

Then, there's no point of $ [0,1] $ where all the functions vanish.
Then, $ f_{1}^{2} + \cdots + f_{n}^{2} $ is a positive function that doesn't vanish at
any point of $ [0,1] $.

\section{Exercise 2.24}

Let $ P,Q $ be distinct maximal ideals of a ring $ R $.

Assume $ P $ contains $ Q^{m} $ for some $ m \geq 1 $.

Since $ P $ is maximal, it's prime, and it's thus radical.

Thus, taking the radicals of $ Q^{m} \subset P $, we get $ Q \subseteq P $,
which is a contradiction.

% I feel like I can use the theory of P-primary modules.



\section{Exercise 7.22(a)}

Recall that irreducible elements correspond to maximal principal ideals.

Let $ r \in R $. If $ r $ is irreducible, we're done.

Otherwise, $ r = r_{1}r_{2} $ for some $ r_{1}, r_{2} \in R $.
If $ r_{1} $ and $ r_{2} $ are irreducible, we're done. Otherwise,
assume without loss of generality that $ r_{1} $ is reducible. Notice that
$ r_{1} \mid r $ if and only if $ (r) \subseteq (r_{1}) $. Thus, this process
has to stop at some point since otherwise we'll produce an infinite ascending
chain of ideals.

\section{Exercise 7.22(b)}

Every ideal $ I $ is finitely generated, so assume $ \sqrt{I} = \langle x_{1}, \ldots, x_{n} \rangle $.

For every $ x_{i} $, $ x_{i}^{k_{i}} \in I $ for some $ k_{i} \geq 1 $.

Let $ k = \sum_{i = 1}^{n}\{k_{i}\} $.

Consider $ \sqrt{I}^{k} $. Then, using the binomial formula, we have the desired result.

\section{Exercise 7.23}

This is an immediate consequence of Proposition 6.7(b).
In particular, considering the contraction of any ascending (descending)
chain of ideals, and then extending back, we get the desired result.



\end{document}
